\documentclass[11pt,a4paper]{article}
\usepackage[utf8]{inputenc}
\usepackage[T1]{fontenc}
\usepackage{lmodern}
\usepackage[margin=1in]{geometry}
\usepackage{graphicx}
\usepackage{booktabs}
\usepackage{longtable}
\usepackage{array}
\usepackage{enumitem}
\usepackage{xcolor}
\usepackage{listings}
\usepackage{hyperref}
\usepackage{fancyhdr}
\usepackage{titlesec}
\usepackage{tocloft}
\usepackage{amsmath}
\usepackage{amssymb}
\usepackage{multicol}
\usepackage{caption}
\usepackage{tabularx}

\definecolor{codeblue}{RGB}{0,102,204}
\definecolor{codegray}{RGB}{100,100,100}
\definecolor{codepurple}{RGB}{153,0,153}
\definecolor{backcolour}{RGB}{248,248,248}
\definecolor{sectionblue}{RGB}{0,51,102}

\lstdefinestyle{matlabstyle}{
  backgroundcolor=\color{backcolour},
  commentstyle=\color{codegray}\itshape,
  keywordstyle=\color{codeblue}\bfseries,
  stringstyle=\color{codepurple},
  basicstyle=\ttfamily\small,
  breaklines=true,
  frame=single,
  rulecolor=\color{codegray!30},
  numbers=left,
  numberstyle=\tiny\color{codegray},
  tabsize=4,
  showstringspaces=false,
  language=Matlab
}
\lstset{style=matlabstyle}

\hypersetup{colorlinks=true,linkcolor=sectionblue,urlcolor=codeblue,citecolor=codepurple}

\pagestyle{fancy}
\fancyhf{}
\fancyhead[L]{\textsc{RiskTool Documentation}}
\fancyhead[R]{\thepage}
\fancyfoot[C]{\small\textcolor{codegray}{Confidential --- For Internal Use Only}}
\renewcommand{\headrulewidth}{0.5pt}

\titleformat{\section}{\Large\bfseries\color{sectionblue}}{\thesection}{1em}{}
\titleformat{\subsection}{\large\bfseries\color{sectionblue!80}}{\thesubsection}{1em}{}
\titleformat{\subsubsection}{\normalsize\bfseries\color{sectionblue!60}}{\thesubsubsection}{1em}{}

\begin{document}

\begin{titlepage}
\centering
\vspace*{2cm}
{\Huge\bfseries\color{sectionblue} RiskTool}\\[0.5cm]
{\LARGE\color{sectionblue!70} Technical Documentation \& User Manual}\\[2cm]
{\Large Portfolio Risk Optimization \& Signal Strategy Framework}\\[1cm]
\rule{\textwidth}{1pt}\\[1.5cm]
{\large\textbf{Modules:}}\\[0.3cm]
{\large Global Asset Allocation (GAA)}\\[0.2cm]
{\large Signal-Driven Equity/Cash Strategy}\\[0.2cm]
{\large Portfolio Risk Optimizer}\\[0.2cm]
{\large Volatility Surface Tooling}\\[2cm]
\vfill
{\large MATLAB R2024b+}\\[0.3cm]
{\large Version 2.0 --- February 2026}\\[0.5cm]
{\small Confidential --- For Internal Use Only}
\end{titlepage}

\tableofcontents
\newpage
\section{System Overview}

\subsection{Introduction}
RiskTool is a comprehensive MATLAB framework for quantitative portfolio management. It provides three integrated modules:
\begin{enumerate}[leftmargin=2em]
  \item \textbf{Portfolio Risk Optimizer} --- Full-featured GUI for risk analysis, optimization, and scenario testing across multiple asset classes.
  \item \textbf{Global Asset Allocation (GAA)} --- Strategic and tactical allocation across 17 asset classes with walk-forward backtesting.
  \item \textbf{Signal Strategy} --- Equity/cash allocation driven by momentum, macro, and composite signals with three strategy modes.
\end{enumerate}
A standalone \textbf{Volatility Surface} utility for SPY options completes the toolkit.

\subsection{Architecture}
The codebase is organized into MATLAB packages with clear separation of concerns:

\begin{center}
\begin{tabularx}{\textwidth}{l l X}
\toprule
\textbf{Layer} & \textbf{Location} & \textbf{Purpose} \\
\midrule
GUI & Root & App Designer GUIs (\texttt{PortfolioRiskOptimizerApp}, \texttt{SignalStrategyApp}, \texttt{SignalBuilderDialog}) \\
Entry Points & Root & Command-line launchers (\texttt{runGAA}, \texttt{run\_*App}) \\
GAA Engine & \texttt{+gaa/} & 17 functions: config, data, CMA, SAA, tactical, backtest, reporting \\
Signal Engine & \texttt{+sigstrat/} & 16 functions: signals, normalization, mapping, 3 strategy modes \\
Utilities & Root & Vol surface (\texttt{GetVolSurface}), test suites \\
\bottomrule
\end{tabularx}
\end{center}

\subsection{Requirements}
\begin{itemize}[leftmargin=2em]
  \item MATLAB R2024b or later (App Designer, \texttt{uifigure}, \texttt{uigridlayout})
  \item \textbf{Optimization Toolbox} --- required for \texttt{fmincon}, \texttt{quadprog}
  \item \textbf{Statistics and Machine Learning Toolbox} --- required for decision trees, cross-validation, \texttt{normcdf}
  \item \textbf{Global Optimization Toolbox} --- optional, for genetic algorithm in robust optimization
  \item Internet access for Stooq and FRED data downloads (cached for 24 hours)
\end{itemize}

\subsection{Quick Start}
\begin{lstlisting}
%% Launch the Portfolio Risk Optimizer GUI
run_PortfolioRiskOptimizerApp

%% Launch the Signal Strategy GUI
run_SignalStrategyApp

%% Run GAA from the command line
results = runGAA('StartDate','20120101', 'Source','Stooq', 'Plot',true);

%% Query implied volatility
iv = GetVolSurface(0.25, 0.95, struct('Type','moneyness','Plot',true));
\end{lstlisting}
\newpage
\section{Portfolio Risk Optimizer}

\subsection{Overview}
The \texttt{PortfolioRiskOptimizerApp} is the main GUI for multi-asset portfolio analysis. It provides:
\begin{itemize}[leftmargin=2em]
  \item Risk metrics: VaR and CVaR via Historical, Parametric, and Monte Carlo methods
  \item Five optimization objectives: Min Variance, Max Sharpe, Min CVaR, Min VaR, Risk Parity (ERC)
  \item Strategic Asset Allocation (SAA) and Tactical Asset Allocation (TAA) with combined weights
  \item Sensitivity analysis ($+1$\% asset shocks) and scenario PnL approximation
  \item Conditional historical bootstrap via user-defined logical expressions
  \item Tracking error analysis and attribution against a benchmark
  \item Performance comparison tab with multiple portfolio configurations
\end{itemize}

\subsection{Launch}
\begin{lstlisting}
run_PortfolioRiskOptimizerApp   %% GUI launcher
app = PortfolioRiskOptimizerApp; %% programmatic access
\end{lstlisting}

\subsection{Tabs and Workflow}
The application is organized into tabs. The typical workflow proceeds left to right:
\begin{enumerate}[leftmargin=2em]
  \item \textbf{Data \& Weights} --- Import or download asset data; define SAA and TAA weights; toggle assets active/inactive.
  \item \textbf{Risk Metrics} --- Compute VaR, CVaR, and contribution analysis at the portfolio level.
  \item \textbf{Optimization} --- Run the selected optimizer subject to constraints (long-only, equity/bond bounds, leverage).
  \item \textbf{Scenario Analysis} --- Define macro scenarios, compute approximate PnL impact, and stress-test the portfolio.
  \item \textbf{Sensitivity} --- Per-asset $+1$\% shock impact on portfolio risk and return.
  \item \textbf{Bootstrap} --- Conditional resampling (e.g., ``VIX $>$ 25 \& SP500\_ret $<$ 0') for tail-risk analysis.
  \item \textbf{Performance} --- Compare strategy vs.\ benchmark cumulative returns, drawdowns, and rolling metrics.
\end{enumerate}

\subsection{Optimization Methods}
\begin{center}
\begin{tabularx}{\textwidth}{l X}
\toprule
\textbf{Method} & \textbf{Description} \\
\midrule
Min Variance & Minimize portfolio variance $w^\top \Sigma w$ subject to constraints \\
Max Sharpe & Maximize $(\mu^\top w - r_f) / \sqrt{w^\top \Sigma w}$ \\
Min CVaR & Minimize expected shortfall at the specified confidence level \\
Min VaR & Minimize Value-at-Risk at the specified confidence level \\
Risk Parity & Equal Risk Contribution --- each asset contributes equally to total portfolio risk \\
\bottomrule
\end{tabularx}
\end{center}
\newpage
\section{Global Asset Allocation (GAA)}

\subsection{Overview}
The GAA module implements a complete strategic and tactical asset allocation framework across 17 asset classes. It can be run from the command line via \texttt{runGAA()} or integrated into the Portfolio Risk Optimizer.

\subsection{Asset Universe}
The default universe (configured in \texttt{gaa.gaaConfig}) spans:
\begin{center}
\begin{tabularx}{\textwidth}{l l X}
\toprule
\textbf{Category} & \textbf{Assets} & \textbf{Tickers (Stooq)} \\
\midrule
Cash \& Bonds & Cash, UST Short, UST Long, TIPS, IG Credit, HY Credit & SHV.US, SHY.US, TLT.US, TIP.US, LQD.US, HYG.US \\
Equities & US Large, US Small, Intl Dev, EM & SPY.US, IWM.US, EFA.US, EEM.US \\
Alternatives & REITs, Commodities, Gold, Hedge Funds & VNQ.US, DJP.US, GLD.US, QAI.US \\
FX & DXY (dollar index) & UUP.US \\
\bottomrule
\end{tabularx}
\end{center}

\subsection{Nine-Step Pipeline}
The \texttt{runGAA()} function executes the following pipeline:

\begin{enumerate}[leftmargin=2em]
  \item \textbf{Configuration} (\texttt{gaa.gaaConfig}) --- Load asset universe, tickers, constraints, scenario definitions, and method selections.
  \item \textbf{Market Data} (\texttt{gaa.downloadData}) --- Download daily prices from Stooq (or Bloomberg/Datastream), compute log returns, align dates. Results cached for 24 hours.
  \item \textbf{Macro Data} (\texttt{gaa.downloadFRED}) --- Download 12 FRED series and compute 14 derived measures (VIX, yield slope, credit spreads, GDP/CPI growth, etc.).
  \item \textbf{Capital Market Assumptions} (\texttt{gaa.computeCMA}) --- Estimate expected returns ($\mu$), volatilities, and the covariance matrix ($\Sigma$) via one of three methods.
  \item \textbf{Strategic Allocation} (\texttt{gaa.computeSAA}) --- Optimize long-term weights using one of five methods, subject to equity/bond constraints.
  \item \textbf{Tactical Signals} (\texttt{gaa.computeTacticalSignals}) --- Compute momentum, risk, and macro signals for tactical overlay.
  \item \textbf{Combine Allocations} (\texttt{gaa.combineAllocations}) --- Blend SAA with tactical signals via Black--Litterman or linear overlay.
  \item \textbf{Walk-Forward Backtest} (\texttt{gaa.backtestStrategy}) --- Historical simulation with rebalancing, transaction costs, and turnover tracking.
  \item \textbf{Analysis \& Reporting} (\texttt{gaa.computePerformance}, \texttt{gaa.plotDashboard}, \texttt{gaa.generateReport}) --- Performance metrics, 6-panel dashboard, and \LaTeX/PDF report.
\end{enumerate}

\subsection{Capital Market Assumptions (CMA)}
\begin{center}
\begin{tabularx}{\textwidth}{l X}
\toprule
\textbf{Method} & \textbf{Description} \\
\midrule
BuildingBlock & Research Affiliates--style forward estimates combining FRED macro measures with equity risk premia, credit spreads, and real rates. \\
Historical & Sample mean and covariance from historical returns, with Ledoit--Wolf shrinkage. \\
BLEquilibrium & Reverse-optimize implied returns from market-cap weights using the Black--Litterman framework. \\
\bottomrule
\end{tabularx}
\end{center}

\subsection{SAA Optimization Methods}
\begin{center}
\begin{tabularx}{\textwidth}{l X}
\toprule
\textbf{Method} & \textbf{Description} \\
\midrule
MinimaxRegret & Robust allocation across 4 growth/inflation scenarios; minimizes worst-case regret. \\
ScenarioWeighted & Scenario probability--weighted optimal portfolios. \\
MaxSharpe & Standard mean--variance optimization for maximum Sharpe ratio. \\
MinVariance & Global minimum variance portfolio. \\
RiskParity & Equal Risk Contribution (ERC) --- each asset contributes equally to total risk. \\
\bottomrule
\end{tabularx}
\end{center}

\subsection{Tactical Overlay}
Tactical signals are computed from three categories:
\begin{itemize}[leftmargin=2em]
  \item \textbf{Momentum} --- Time-series (vol-scaled cumulative return) and cross-sectional (demeaned) momentum per asset.
  \item \textbf{Risk} --- VIX level, credit spread regime, and volatility regime indicators.
  \item \textbf{Macro} --- Yield curve slope, growth, and inflation momentum.
\end{itemize}
Signals are combined into a composite $\in [-1,+1]$ per asset. Two blending modes are available:
\begin{itemize}[leftmargin=2em]
  \item \textbf{Black--Litterman} --- Convert signals to views and blend with SAA via the BL posterior.
  \item \textbf{Linear} --- Scale signals by the tactical budget and apply as additive deltas to SAA weights.
\end{itemize}

\subsection{Covariance Estimation}
Robust covariance estimation is provided by \texttt{gaa.robustCovariance}:
\begin{itemize}[leftmargin=2em]
  \item \textbf{Ledoit--Wolf Shrinkage} --- Shrinks sample covariance toward a scaled identity target, with analytically optimal intensity $\delta \in [0,1]$.
  \item \textbf{Marc\v{e}nko--Pastur Denoising} --- Optional eigenvalue clipping that replaces noise eigenvalues with their average, improving out-of-sample stability.
\end{itemize}

\subsection{Risk Decomposition}
The function \texttt{gaa.computeRiskDecomp} provides:
\begin{itemize}[leftmargin=2em]
  \item \textbf{MCR} --- Marginal Contribution to Risk: $\partial\sigma_p / \partial w_i$.
  \item \textbf{CR} --- Contribution to Risk: $w_i \times \text{MCR}_i$.
  \item \textbf{PCT} --- Percentage risk contribution: $\text{CR}_i / \sigma_p$.
  \item \textbf{HHI} --- Herfindahl--Hirschman Index of risk concentration.
  \item \textbf{ENB} --- Effective Number of Bets: $1/\text{HHI}$.
  \item \textbf{DR} --- Diversification Ratio: $w^\top\sigma / \sigma_p$.
\end{itemize}

\subsection{Command-Line Usage}
\begin{lstlisting}
%% Full pipeline with defaults
results = runGAA();

%% Custom configuration
results = runGAA('StartDate', '20150101', ...
    'Source',    'Stooq', ...
    'CMAMethod', 'BuildingBlock', ...
    'SAAMethod', 'MinimaxRegret', ...
    'CombineMethod', 'BL', ...
    'TacticalBudget', 0.20, ...
    'Plot', true);
\end{lstlisting}
\newpage
\section{Signal Strategy}

\subsection{Overview}
The Signal Strategy module implements an equity/cash allocation system driven by quantitative signals. It supports three distinct strategy modes, each with its own optimization approach.

\subsection{Launch}
\begin{lstlisting}
run_SignalStrategyApp   %% GUI
\end{lstlisting}

\subsection{Data Pipeline}
\begin{enumerate}[leftmargin=2em]
  \item \textbf{Market Data} (\texttt{sigstrat.downloadMarketData}) --- S\&P 500 and T-Bill daily returns from Stooq, with SPY.US fallback and 24h caching.
  \item \textbf{Macro Data} (\texttt{sigstrat.downloadMacroData}) --- 12 FRED series with 14 derived measures; shares pipeline with GAA where available.
  \item \textbf{Signal Generation} --- Three sources:
  \begin{itemize}
    \item Demo signals (\texttt{sigstrat.generateDemoSignals}) --- 5 quick signals for exploration
    \item Extended signals (\texttt{sigstrat.generateExtendedSignals}) --- 12 production signals across fast/medium/slow categories
    \item Custom signals (\texttt{SignalBuilderDialog}) --- Interactive builder with 6 signal types
  \end{itemize}
\end{enumerate}

\subsection{Signal Types}
The \texttt{sigstrat.buildSignal} function supports six signal types, all producing output bounded in $[-1, +1]$:

\subsubsection{Momentum}
Vol-adjusted price momentum with expanding z-score and configurable activation function.

\begin{center}
\begin{tabularx}{\textwidth}{l l l X}
\toprule
\textbf{Parameter} & \textbf{Default} & \textbf{Range} & \textbf{Description} \\
\midrule
\texttt{lookback} & 126 & 21--504 & Lookback window in trading days \\
\texttt{minHistory} & 252 & 42--1008 & Minimum observations for expanding z-score \\
\texttt{volHalflife} & 32 & 10--120 & EWMA volatility halflife (days) \\
\texttt{skipDays} & 0 & 0--21 & Skip recent days to avoid short-term reversal \\
\texttt{volFloor} & 0 & $\geq 0$ & Min vol estimate; 0 = auto (5th percentile) \\
\texttt{activation} & tanh & --- & Activation function (see below) \\
\texttt{tanhScale} & 2 & 0.1--10 & Scaling divisor for activation \\
\bottomrule
\end{tabularx}
\end{center}

The momentum signal at time $t$ is:
\begin{equation}
\text{mom}_t = \frac{\sum_{i=t-L-S+1}^{t-S} r_i}{\hat{\sigma}_t^{\text{EWMA}} \cdot \sqrt{L}}
\end{equation}
where $L$ is \texttt{lookback}, $S$ is \texttt{skipDays}, and $\hat{\sigma}_t^{\text{EWMA}}$ is the EWMA volatility.

\subsubsection{EWMAC (Exponentially Weighted Moving Average Crossover)}
EMA fast/slow crossover trend signal, following Levine--Pedersen (2016). Equivalent to TSMOM but smoother with lower turnover.

\begin{center}
\begin{tabularx}{\textwidth}{l l l X}
\toprule
\textbf{Parameter} & \textbf{Default} & \textbf{Range} & \textbf{Description} \\
\midrule
\texttt{fastSpan} & 16 & 2--128 & Fast EMA span \\
\texttt{slowSpan} & 64 & 8--512 & Slow EMA span \\
\texttt{volHalflife} & 32 & 10--120 & EWMA vol halflife \\
\texttt{volFloor} & 0 & $\geq 0$ & Min vol estimate \\
\texttt{activation} & tanh & --- & Activation function \\
\texttt{tanhScale} & 2 & 0.1--10 & Activation scale \\
\bottomrule
\end{tabularx}
\end{center}

The raw EWMAC signal is:
\begin{equation}
\text{EWMAC}_t = \frac{\text{EMA}_{\text{fast}}(P_t) - \text{EMA}_{\text{slow}}(P_t)}{\hat{\sigma}_t^{\text{EWMA}} \cdot P_t}
\end{equation}

\subsubsection{Ensemble (Multi-Lookback Momentum)}
Combines momentum signals across multiple lookback windows. This is the single most impactful enhancement per practitioners (Moskowitz--Ooi--Pedersen 2012, Baltas--Kosowski 2020).

\begin{center}
\begin{tabularx}{\textwidth}{l l l X}
\toprule
\textbf{Parameter} & \textbf{Default} & \textbf{Range} & \textbf{Description} \\
\midrule
\texttt{lookbacks} & {[}21 63 126 252{]} & --- & Vector of lookback windows \\
\texttt{method} & equal & --- & \texttt{equal} (equal weight) or \texttt{riskparity} ($1/\sigma$ weighting) \\
\texttt{volHalflife} & 32 & 10--120 & EWMA vol halflife \\
\texttt{skipDays} & 0 & 0--21 & Skip recent days \\
\texttt{volFloor} & 0 & $\geq 0$ & Min vol estimate \\
\texttt{activation} & tanh & --- & Activation function \\
\bottomrule
\end{tabularx}
\end{center}

\subsubsection{MacroZ}
Expanding robust z-score of a macro variable (e.g., VIX, HY\_Spread). Uses median/MAD for robustness to outliers.

\begin{center}
\begin{tabularx}{\textwidth}{l l l X}
\toprule
\textbf{Parameter} & \textbf{Default} & \textbf{Range} & \textbf{Description} \\
\midrule
\texttt{variable} & VIX & --- & Macro column name \\
\texttt{minWindow} & 126 & 20--504 & Min observations for z-score \\
\texttt{negate} & false & --- & Flip sign (e.g., lower VIX = bullish) \\
\texttt{tanhScale} & 2 & 0.1--10 & Scaling divisor \\
\bottomrule
\end{tabularx}
\end{center}

\subsubsection{MeanReversion}
Distance from rolling moving average, z-scored and sign-flipped. Price above MA $\rightarrow$ expect reversion down $\rightarrow$ bearish signal.

\subsubsection{Composite}
Weighted average of multiple macro z-scores. User defines the variables, weights, and per-variable sign flips.

\subsection{EWMA Volatility Estimator}
All momentum-family signals (Momentum, EWMAC, Ensemble) use EWMA volatility instead of equal-weight standard deviation:
\begin{equation}
\hat{\sigma}^2_t = \lambda \hat{\sigma}^2_{t-1} + (1 - \lambda) r_t^2, \quad \lambda = \exp\left(-\frac{\ln 2}{h}\right)
\end{equation}
where $h$ is the halflife in days (default 32). This places more weight on recent observations, producing a smoother and more responsive vol estimate.

\subsection{Volatility Floor}
To prevent extreme leverage in low-volatility regimes, a vol floor is applied:
\begin{equation}
\hat{\sigma}_t = \max(\hat{\sigma}_t^{\text{EWMA}},\; \text{floor})
\end{equation}
When \texttt{volFloor} = 0 (auto), the floor is the 5th percentile of the trailing 504-day EWMA vol series. When $> 0$, it is used directly.

\subsection{Activation Functions}
Three activation functions map the z-scored signal to $[-1, +1]$:

\begin{center}
\begin{tabularx}{\textwidth}{l l X}
\toprule
\textbf{Name} & \textbf{Formula} & \textbf{Behavior} \\
\midrule
\texttt{tanh} & $\tanh(z/s)$ & Standard S-curve, saturates at $\pm 1$ (default) \\
\texttt{sigmoid} & $2\,\Phi(z/s) - 1$ & Normal CDF scaled to $[-1,+1]$; lighter tails than tanh \\
\texttt{revertSigmoid} & $c \cdot (z/s) \cdot e^{-z^2/(2s^2)}$ & Dao et al.\ (2021): position rises then falls for extreme signals, providing crash protection \\
\bottomrule
\end{tabularx}
\end{center}
where $s$ = \texttt{tanhScale}, $\Phi$ = standard normal CDF, and $c = (1 + 2/s^2)^{3/4}$.

\subsection{Extended Signal Set (12 Signals)}
The function \texttt{sigstrat.generateExtendedSignals} produces 12 production signals:

\begin{center}
\begin{tabularx}{\textwidth}{l l l X}
\toprule
\textbf{\#} & \textbf{Name} & \textbf{Speed} & \textbf{Description} \\
\midrule
1 & SPX\_Mom\_1m & Fast & 21d equity momentum \\
2 & SPX\_Mom\_3m & Fast & 63d equity momentum \\
3 & VIX\_Fast & Fast & Negated 63d VIX z-score \\
4 & CreditFast & Fast & Negated 63d HY spread z-score \\
5 & SPX\_Mom\_6m & Medium & 126d equity momentum \\
6 & VIX\_Med & Medium & Negated 126d VIX z-score \\
7 & YieldSlope\_Med & Medium & 126d yield slope z-score \\
8 & FinConditions & Medium & VIX + credit + yield composite \\
9 & SPX\_Mom\_12m & Slow & 252d equity momentum \\
10 & VIX\_Slow & Slow & Negated 252d VIX z-score \\
11 & MacroTrend & Slow & GDP + retail + industrial composite \\
12 & CreditSlow & Slow & Negated 252d HY spread z-score \\
\bottomrule
\end{tabularx}
\end{center}

\subsection{Strategy Mode 1: Composite Signal}
The \texttt{sigstrat.optimizeComposite} function:
\begin{enumerate}[leftmargin=2em]
  \item Computes a weighted average of normalized signals: $C_t = \sum_i w_i \cdot s_{i,t}$
  \item Maps $C_t$ to equity weight via step-function thresholds
  \item Optimizes weights and thresholds via grid search (in-sample Sharpe + return + drawdown penalty)
  \item Reports in-sample and out-of-sample performance
\end{enumerate}

\subsection{Strategy Mode 2: Risk Budget}
The \texttt{sigstrat.riskBudgetStrategy} function:
\begin{enumerate}[leftmargin=2em]
  \item Allocates a risk budget to each signal (default: equal)
  \item Maps each signal to equity weight via a configurable translation function (Linear, Threshold, Sigmoid, Quantile, Piecewise, Spline, or Power)
  \item Combines per-signal allocations weighted by budget
  \item Supports optional macro conditioning
\end{enumerate}

\subsection{Strategy Mode 3: Decision Tree}
The \texttt{sigstrat.decisionTreeStrategy} function:
\begin{enumerate}[leftmargin=2em]
  \item \textbf{Rule-Based} --- User-defined rules (variable, operator, threshold $\rightarrow$ equity weight)
  \item \textbf{Auto-Optimize} --- CART regression tree auto-fitted to predict equity weight $\in [0,1]$
  \item Walk-forward cross-validation with expanding windows
  \item Variable importance metrics from the fitted tree
\end{enumerate}

\subsection{Signal Normalization}
The \texttt{sigstrat.normalizeSignals} function supports 5 methods:
\begin{center}
\begin{tabularx}{\textwidth}{l X}
\toprule
\textbf{Method} & \textbf{Description} \\
\midrule
RobustZ & Expanding robust z-score (median/MAD) $\rightarrow$ tanh (default) \\
StandardZ & Expanding standard z-score $\rightarrow$ tanh \\
RollingZ & Rolling-window z-score $\rightarrow$ tanh \\
Percentile & Expanding percentile rank mapped to $[-1, +1]$ \\
MinMax & Expanding min-max scaling to $[-1, +1]$ \\
\bottomrule
\end{tabularx}
\end{center}

\subsection{Signal-to-Weight Mapping}
The \texttt{sigstrat.mapSignalToWeight} function maps signals from $[-1,+1]$ to equity weights $[0,1]$:
\begin{center}
\begin{tabularx}{\textwidth}{l X}
\toprule
\textbf{Method} & \textbf{Description} \\
\midrule
Step & Threshold-based discrete levels \\
Linear & $(\text{signal} + 1) / 2$ \\
Sigmoid & Smooth sigmoid mapping \\
PiecewiseLinear & User-defined breakpoints \\
Spline & Smooth spline interpolation \\
Power & Power law transformation \\
\bottomrule
\end{tabularx}
\end{center}

\subsection{Backtesting}
The \texttt{sigstrat.backtestEquityCash} function performs walk-forward simulation:
\begin{itemize}[leftmargin=2em]
  \item Configurable rebalancing frequency (daily, weekly, monthly)
  \item Transaction cost modeling (basis points per trade)
  \item Weight drift tracking between rebalance dates
  \item Benchmark comparisons (60/40, 100\% equity)
\end{itemize}

\subsection{Stress Testing}
The \texttt{sigstrat.stressTestStrategy} function evaluates performance during 6 historical crisis episodes:
\begin{center}
\begin{tabularx}{\textwidth}{l l X}
\toprule
\textbf{Episode} & \textbf{Period} & \textbf{Event} \\
\midrule
GFC & Oct 2007 -- Mar 2009 & Global Financial Crisis \\
Euro Crisis & Jul 2011 -- Oct 2011 & European sovereign debt \\
Taper Tantrum & May 2013 -- Sep 2013 & Fed taper announcement \\
2018 Q4 & Oct 2018 -- Dec 2018 & Rate hike selloff \\
COVID & Feb 2020 -- Mar 2020 & Pandemic crash \\
Rate Shock & Jan 2022 -- Oct 2022 & Inflation and rate hiking cycle \\
\bottomrule
\end{tabularx}
\end{center}
\newpage
\section{Signal Builder Dialog}

\subsection{Overview}
The \texttt{SignalBuilderDialog} is a modal UI launched from the Signal Strategy App for building individual custom signals interactively.

\subsection{Workflow}
\begin{enumerate}[leftmargin=2em]
  \item Select signal type from the dropdown (6 types available)
  \item Configure parameters in the dynamic panel (parameters change per type)
  \item Enter a unique signal name
  \item Click \textbf{Preview} to visualize the signal and inspect statistics
  \item Click \textbf{Add Signal} to accept, or \textbf{Cancel} to discard
\end{enumerate}

\subsection{Signal Type Panels}
Each signal type presents a tailored parameter panel:
\begin{itemize}[leftmargin=2em]
  \item \textbf{Momentum} --- Lookback, Min History, Vol Halflife, Skip Days, Vol Floor, Activation, Tanh Scale
  \item \textbf{EWMAC} --- Fast Span, Slow Span, Vol Halflife, Vol Floor, Activation, Tanh Scale
  \item \textbf{Ensemble} --- Lookbacks (comma-separated text), Method (Equal Weight / Risk Parity), Vol Halflife, Skip Days, Vol Floor, Activation, Tanh Scale
  \item \textbf{MacroZ} --- Variable dropdown, Min Window, Negate checkbox, Tanh Scale
  \item \textbf{MeanReversion} --- MA Lookback, Min History, Tanh Scale
  \item \textbf{Composite} --- Editable table (Variable, Weight, Negate per row), Min Window, Tanh Scale
\end{itemize}
\newpage
\section{Volatility Surface Utility}

\subsection{Overview}
The \texttt{GetVolSurface} function downloads SPY options, calibrates an SVI (Stochastic Volatility Inspired) surface, and queries implied volatility at any $(T, K)$ point.

\subsection{Usage}
\begin{lstlisting}
%% Query IV at 3-month, 95% moneyness
iv = GetVolSurface(0.25, 0.95, struct('Type','moneyness'));

%% Query IV at 1-year, strike = 500
iv = GetVolSurface(1.0, 500, struct('Type','strike'));

%% Force refresh and plot 3D surface
iv = GetVolSurface(0.5, 1.0, struct('Type','moneyness', ...
    'Refresh',true, 'Plot',true));
\end{lstlisting}

\subsection{SVI Model}
The SVI parameterization models total implied variance as a function of log-moneyness $k$:
\begin{equation}
w(k) = a + b\left(\rho(k - m) + \sqrt{(k - m)^2 + \sigma^2}\right)
\end{equation}
Parameters $(a, b, \rho, m, \sigma)$ are calibrated per expiry via least-squares fit to market quotes. The function uses persistent caching (\texttt{VolSurfaceCache.mat}) and a 24-hour refresh policy.
\newpage
\section{GAA Function Reference}

\begin{longtable}{p{3.8cm} p{11.2cm}}
\toprule
\textbf{Function} & \textbf{Description} \\
\midrule
\endhead
\texttt{gaaConfig} & Central configuration: asset universe (17 assets), tickers, FRED series, scenario definitions, optimization constraints, and method selections. Returns a struct used by all downstream functions. \\
& {\small\texttt{cfg = gaa.gaaConfig()}} \\
\midrule
\texttt{downloadData} & Download daily prices from Stooq/Bloomberg/Datastream. Computes log returns, aligns dates across assets. Caches results for 24 hours. \\
& {\small\texttt{[R, prices, dates, names] = gaa.downloadData(cfg)}} \\
\midrule
\texttt{downloadFRED} & Download 12 FRED series and compute 14 derived macro measures (VIX, yield slope, credit spreads, GDP/CPI growth, DXY, NFP, etc.). \\
& {\small\texttt{macro = gaa.downloadFRED(cfg)}} \\
\midrule
\texttt{computeCMA} & Estimate Capital Market Assumptions: expected returns, volatilities, and covariance matrix via BuildingBlock, Historical, or BLEquilibrium methods. \\
& {\small\texttt{cma = gaa.computeCMA(R, cfg, macro)}} \\
\midrule
\texttt{computeSAA} & Compute Strategic Asset Allocation via MinimaxRegret, ScenarioWeighted, MaxSharpe, MinVariance, or RiskParity. \\
& {\small\texttt{saa = gaa.computeSAA(cma, cfg)}} \\
\midrule
\texttt{computeTacticalSignals} & Compute tactical signals from momentum, risk, and macro categories. Returns composite signal per asset in [-1,+1]. \\
& {\small\texttt{signals = gaa.computeTacticalSignals(R, macro, cfg)}} \\
\midrule
\texttt{combineAllocations} & Blend SAA with tactical signals via Black-Litterman or linear overlay. Enforces tactical budget and per-asset deviation constraints. \\
& {\small\texttt{alloc = gaa.combineAllocations(saa, signals, cma, cfg)}} \\
\midrule
\texttt{backtestStrategy} & Walk-forward historical backtest with rebalancing, transaction costs, turnover tracking, and SAA benchmark. \\
& {\small\texttt{bt = gaa.backtestStrategy(alloc, R, cfg)}} \\
\midrule
\texttt{computePerformance} & Comprehensive performance analytics: annualized metrics, drawdown, rolling statistics, monthly returns, attribution. \\
& {\small\texttt{perf = gaa.computePerformance(bt)}} \\
\midrule
\texttt{computeRiskDecomp} & Risk decomposition: MCR, CR, percentage contributions, HHI, ENB, and Diversification Ratio. \\
& {\small\texttt{rd = gaa.computeRiskDecomp(w, Sigma)}} \\
\midrule
\texttt{optimizePortfolio} & Unified optimizer: MaxSharpe, MinVariance, RiskParity, MinCVaR, TargetVol, Resampled. Long-only and constrained. \\
& {\small\texttt{w = gaa.optimizePortfolio(mu, Sigma, opts)}} \\
\midrule
\texttt{blackLitterman} & Black-Litterman posterior returns and covariance from views. \\
& {\small\texttt{[muBL, SigBL] = gaa.blackLitterman(Sigma, wMkt, P, Q, opts)}} \\
\midrule
\texttt{robustCovariance} & Ledoit-Wolf shrinkage with optional Marcenko-Pastur eigenvalue denoising. \\
& {\small\texttt{[S, info] = gaa.robustCovariance(R, opts)}} \\
\midrule
\texttt{ledoitWolfShrink} & Ledoit-Wolf (2004) shrinkage toward scaled identity. \\
& {\small\texttt{[S, delta] = gaa.ledoitWolfShrink(R)}} \\
\midrule
\texttt{plotDashboard} & 6-panel dashboard: CMA bars, signal heatmap, allocation area, risk waterfall, cumulative performance, drawdown. \\
& {\small\texttt{fig = gaa.plotDashboard(results)}} \\
\midrule
\texttt{generateReport} & LaTeX/PDF report generation with embedded dashboard figure. \\
& {\small\texttt{gaa.generateReport(results)}} \\
\bottomrule
\end{longtable}
\newpage
\section{Signal Strategy Function Reference}

\begin{longtable}{p{3.8cm} p{11.2cm}}
\toprule
\textbf{Function} & \textbf{Description} \\
\midrule
\endhead
\texttt{downloadMarketData} & Download S\\&P 500 and T-Bill daily data from Stooq. Returns struct with dates, returns, and prices. 24h cache. \\
& {\small\texttt{mkt = sigstrat.downloadMarketData()}} \\
\midrule
\texttt{downloadMacroData} & Download FRED macro series with 14 derived measures. Falls back to embedded pipeline if gaa.downloadFRED unavailable. \\
& {\small\texttt{macro = sigstrat.downloadMacroData(dates)}} \\
\midrule
\texttt{generateDemoSignals} & Generate 5 demo signals: yield slope, VIX, credit spread, momentum, macro composite. \\
& {\small\texttt{[data, names] = sigstrat.generateDemoSignals(macro, mkt)}} \\
\midrule
\texttt{generateExtendedSignals} & Generate 12 production signals across fast/medium/slow categories. Uses EWMA vol. \\
& {\small\texttt{[data, names] = sigstrat.generateExtendedSignals(macro, mkt)}} \\
\midrule
\texttt{buildSignal} & Build custom signal: Momentum, EWMAC, Ensemble, MacroZ, MeanReversion, or Composite. \\
& {\small\texttt{[sig, meta] = sigstrat.buildSignal(type, params, mkt, macro)}} \\
\midrule
\texttt{normalizeSignals} & Normalize signals to [-1,+1] via RobustZ, StandardZ, RollingZ, Percentile, or MinMax. \\
& {\small\texttt{normed = sigstrat.normalizeSignals(data, method)}} \\
\midrule
\texttt{mapSignalToWeight} & Map signal [-1,+1] to equity weight [0,1] via Step, Linear, Sigmoid, Piecewise, Spline, or Power. \\
& {\small\texttt{w = sigstrat.mapSignalToWeight(sig, method, opts)}} \\
\midrule
\texttt{backtestEquityCash} & Walk-forward equity/cash backtest with rebalancing, costs, and benchmark comparisons. \\
& {\small\texttt{bt = sigstrat.backtestEquityCash(w, mkt, opts)}} \\
\midrule
\texttt{computeStrategyPerf} & Performance metrics for strategy, 60/40, and 100\\% equity. Rolling analytics and monthly heatmap. \\
& {\small\texttt{perf = sigstrat.computeStrategyPerf(bt)}} \\
\midrule
\texttt{stressTestStrategy} & Evaluate strategy during 6 historical crisis episodes. \\
& {\small\texttt{st = sigstrat.stressTestStrategy(bt)}} \\
\midrule
\texttt{optimizeComposite} & Mode 1: grid-search optimal weights and thresholds for composite signal strategy. \\
& {\small\texttt{res = sigstrat.optimizeComposite(signals, mkt, opts)}} \\
\midrule
\texttt{riskBudgetStrategy} & Mode 2: per-signal risk budget allocation with configurable translation functions. \\
& {\small\texttt{res = sigstrat.riskBudgetStrategy(signals, mkt, opts)}} \\
\midrule
\texttt{robustOptimize} & K-fold CV robust optimization with L2 regularization and turnover penalty. \\
& {\small\texttt{res = sigstrat.robustOptimize(signals, mkt, opts)}} \\
\midrule
\texttt{decisionTreeStrategy} & Mode 3: rule-based or CART auto-optimized decision tree for weight prediction. \\
& {\small\texttt{res = sigstrat.decisionTreeStrategy(signals, mkt, opts)}} \\
\bottomrule
\end{longtable}
\newpage
\section{Data Structures Reference}

\subsection{Market Data (\texttt{mkt} struct)}
\begin{center}
\begin{tabularx}{\textwidth}{l l X}
\toprule
\textbf{Field} & \textbf{Type} & \textbf{Description} \\
\midrule
\texttt{dates} & datetime & Price dates \\
\texttt{retDates} & datetime & Return dates (\texttt{dates(2:end)}) \\
\texttt{spxRet} & Tx1 double & Daily S\&P 500 log returns \\
\texttt{cashRet} & Tx1 double & Daily T-Bill returns \\
\texttt{spxPrice} & Tx1 double & S\&P 500 price level \\
\bottomrule
\end{tabularx}
\end{center}

\subsection{Macro Data (\texttt{macro} table)}
A MATLAB table with \texttt{Date} column plus derived measures:
\begin{multicols}{3}
\begin{itemize}[leftmargin=1em,itemsep=0pt]
  \item VIX
  \item FedFunds
  \item TenYr
  \item TwoYr
  \item YieldSlope
  \item HY\_Spread
  \item IG\_Spread
  \item GDP\_YoY
  \item CPI\_YoY
  \item CPI\_MoM
  \item RetailSales\_YoY
  \item IndProd\_YoY
  \item NFP\_MoM
  \item DXY\_3m
\end{itemize}
\end{multicols}

\subsection{CMA Struct}
\begin{center}
\begin{tabularx}{\textwidth}{l l X}
\toprule
\textbf{Field} & \textbf{Type} & \textbf{Description} \\
\midrule
\texttt{mu} & Nx1 & Expected annualized returns \\
\texttt{vol} & Nx1 & Expected annualized volatilities \\
\texttt{Sigma} & NxN & Covariance matrix \\
\texttt{Corr} & NxN & Correlation matrix \\
\texttt{method} & char & Method used (BuildingBlock, Historical, BLEquilibrium) \\
\bottomrule
\end{tabularx}
\end{center}

\subsection{Performance Struct}
\begin{center}
\begin{tabularx}{\textwidth}{l l X}
\toprule
\textbf{Field} & \textbf{Type} & \textbf{Description} \\
\midrule
\texttt{annReturn} & scalar & Annualized return \\
\texttt{annVol} & scalar & Annualized volatility \\
\texttt{sharpe} & scalar & Sharpe ratio \\
\texttt{sortino} & scalar & Sortino ratio \\
\texttt{calmar} & scalar & Calmar ratio (return / max drawdown) \\
\texttt{maxDD} & scalar & Maximum drawdown \\
\texttt{hitRate} & scalar & Fraction of positive return days \\
\texttt{skew} & scalar & Return skewness \\
\texttt{kurt} & scalar & Return kurtosis \\
\bottomrule
\end{tabularx}
\end{center}
\newpage
\section{Testing}

\subsection{Test Suites}
Two test suites validate the framework:

\subsubsection{GAA End-to-End Test (\texttt{testGAA\_synthetic.m})}
Uses synthetic data (10 years, 17 assets, 3-factor correlation structure) to test every module in the GAA pipeline without requiring internet access. Validates:
\begin{itemize}[leftmargin=2em]
  \item Data download and alignment
  \item CMA computation (all 3 methods)
  \item SAA optimization (all 5 methods)
  \item Tactical signal generation
  \item Allocation blending
  \item Backtest execution
  \item Performance computation
\end{itemize}

\subsubsection{Momentum Upgrades Test (\texttt{test\_momentum\_upgrades.m})}
Unit tests for the enhanced momentum signal system. 12 test cases covering:
\begin{itemize}[leftmargin=2em]
  \item Momentum default parameters and valid output range
  \item EWMA vol produces different signals with different halflife values
  \item SkipDays produces shifted signals
  \item All 3 activation functions produce valid $[-1,+1]$ output
  \item EWMAC with default and custom spans
  \item Ensemble with equal-weight and risk-parity methods
  \item Vol floor prevents extreme values
  \item Extended signals produce $T \times 12$ output
  \item Backward compatibility of MacroZ
\end{itemize}

\subsection{Running Tests}
\begin{lstlisting}
%% Run all tests
runtests('testGAA_synthetic');
runtests('test_momentum_upgrades');

%% Run with verbose output
results = runtests('test_momentum_upgrades', 'Verbosity', 3);
disp(results);
\end{lstlisting}
\newpage
\section{Caching and Data Management}

\subsection{Data Caching}
All data download functions implement a 24-hour disk cache to minimize API calls:
\begin{center}
\begin{tabularx}{\textwidth}{l l X}
\toprule
\textbf{Module} & \textbf{Cache Location} & \textbf{Contents} \\
\midrule
\texttt{gaa.downloadData} & \texttt{tempdir/gaa\_cache/} & Market prices, returns, dates \\
\texttt{gaa.downloadFRED} & \texttt{tempdir/gaa\_cache/} & FRED macro series \\
\texttt{sigstrat.downloadMarketData} & \texttt{tempdir/sigstrat\_cache/} & S\&P 500, T-Bill data \\
\texttt{GetVolSurface} & \texttt{VolSurfaceCache.mat} & SVI parameters per expiry \\
\bottomrule
\end{tabularx}
\end{center}

To force a data refresh, delete the cache directory or pass the \texttt{Refresh} option where available.

\subsection{Output Directory}
The GAA module writes reports and figures to \texttt{gaa\_output/}:
\begin{itemize}[leftmargin=2em]
  \item \texttt{GAA\_Report.pdf} --- Auto-generated \LaTeX/PDF strategy report
  \item Dashboard PNG exports
\end{itemize}
\newpage
\section{Academic References}

The signal construction and portfolio methodology draws on the following research:
\begin{itemize}[leftmargin=2em]
  \item Moskowitz, T., Ooi, Y.H., \& Pedersen, L.H. (2012). ``Time series momentum.' \textit{Journal of Financial Economics}, 104(2), 228--250.
  \item Baltas, N. \& Kosowski, R. (2013/2020). ``Momentum strategies in futures markets and trend-following funds.' Working paper / \textit{Journal of Finance and Data Science}.
  \item Levine, A. \& Pedersen, L.H. (2016). ``Which trend is your friend?' \textit{Financial Analysts Journal}, 72(3), 51--66.
  \item Dao, T.-L., Nguyen, T.-T., Deremble, C., Lem\'erle, Y., Bouchaud, J.-P., \& Potters, M. (2021). ``Tail protection for long investors: Trend convexity at work.' Working paper.
  \item Novy-Marx, R. (2012). ``Is momentum really momentum?' \textit{Journal of Financial Economics}, 103(3), 429--453.
  \item Black, F. \& Litterman, R. (1992). ``Global portfolio optimization.' \textit{Financial Analysts Journal}, 48(5), 28--43.
  \item Ledoit, O. \& Wolf, M. (2004). ``A well-conditioned estimator for large-dimensional covariance matrices.' \textit{Journal of Multivariate Analysis}, 88(2), 365--411.
  \item Michaud, R.O. (1998). \textit{Efficient Asset Management}. Harvard Business School Press.
  \item Carver, R. (2015). \textit{Systematic Trading}. Harriman House.
  \item Gatheral, J. (2004). ``A parsimonious arbitrage-free implied volatility parameterization with application to the valuation of volatility derivatives.' Presentation, Global Derivatives.
\end{itemize}

\end{document}
